\chapter{Methodology and System Design}
\label{chap:methodology}
% How the project works and how it was built: the most reproducible chapter.
% Covers CRISP-DM Data Understanding, Data Preparation, Modeling (MO.1, MO.3)
% and Deployment planning (DE.1, DE.2). MO.2 and MO.4 are in Chapter 4.

\section{Introduction}

\section{Development Methodology: CRISP-DM}
\label{sec:crispdm}
The project follows the Cross-Industry Standard Process for Data Mining
(CRISP-DM)~\cite{chapman2000crisp,wirth2000crisp}, which organizes a data
project into six phases (business understanding, data understanding, data
preparation, modeling, evaluation and deployment), each broken down into
generic tasks with defined outputs. The phases are not executed strictly in
sequence: the process explicitly allows returning to an earlier phase when a
later one reveals a problem, as shown in Figure~\ref{fig:crispdm-cycle}.

% Then, about THIS project:
% - why CRISP-DM rather than an alternative such as the KDD process
%   \cite{fayyad1996kdd} (cite only verified sources);
% - how it was adapted (e.g. which tasks were light or not applicable);
% - the iterations that ACTUALLY happened between phases, with evidence
%   (experiment logs, git history). Do not invent loops.
% - a pointer to Table~\ref{tab:crispdm-traceability}, which maps every
%   generic task to the section of this report that addresses it.

\begin{figure}[htbp]
  \centering
  % CRISP-DM process model, redrawn after Chapman et al. (2000).
% Needs: \usepackage{tikz} and \usetikzlibrary{arrows.meta,shapes.geometric}.
\begin{tikzpicture}[
    font=\small,
    phase/.style={draw, thick, rounded corners=3pt, fill=gray!8, align=center,
                  text width=2.4cm, minimum height=1.1cm},
    flow/.style={-{Stealth[length=2.2mm]}, thick},
  ]
  \node[draw, thick, cylinder, shape border rotate=90, aspect=0.25,
        fill=gray!20, minimum width=1.6cm, minimum height=1.5cm] (data) at (0,0) {Data};

  \node[phase] (bu) at (-3.2, 2.7) {Business\\understanding};
  \node[phase] (du) at ( 3.2, 2.7) {Data\\understanding};
  \node[phase] (dp) at ( 4.4, 0.0) {Data\\preparation};
  \node[phase] (mo) at ( 3.2,-2.7) {Modeling};
  \node[phase] (ev) at ( 0.0,-3.8) {Evaluation};
  \node[phase] (de) at (-4.4,-1.2) {Deployment};

  % Transitions of the reference model; paired arrows mean frequent back-and-forth.
  \draw[flow] (bu) to[bend left=12] (du);
  \draw[flow] (du) to[bend left=12] (bu);
  \draw[flow] (du) -- (dp);
  \draw[flow] (dp) to[bend left=12] (mo);
  \draw[flow] (mo) to[bend left=12] (dp);
  \draw[flow] (mo) -- (ev);
  \draw[flow] (ev) -- (bu);
  \draw[flow] (ev) -- (de);

  % Outer cycle: data projects iterate.
  \draw[flow, gray!70, line width=1.2pt]
    (100:6.3) arc[start angle=100, end angle=-240, radius=6.3];
\end{tikzpicture}

  \caption[The CRISP-DM process model]{The six phases of the CRISP-DM process
    model and the main transitions between them; the outer arrow denotes the
    cyclic nature of data projects (redrawn after~\cite{chapman2000crisp}).}
  \label{fig:crispdm-cycle}
\end{figure}

% Traceability of the 24 CRISP-DM generic tasks to this report.
% The last column fills itself from the \crispdm{ID} markers: do not type
% section numbers by hand. DE.3 and DE.4 live in unnumbered chapters, so they
% link by name instead of by section number.
\begingroup
\setstretch{1.05}\small
\begin{longtable}{@{}p{0.23\textwidth}p{0.07\textwidth}p{0.38\textwidth}p{0.22\textwidth}@{}}
\caption[Traceability of the CRISP-DM generic tasks to this report]{Traceability
  of the CRISP-DM generic tasks~\cite{chapman2000crisp} to the sections of this
  report.}
\label{tab:crispdm-traceability}\\
\toprule
\textbf{Phase} & \textbf{ID} & \textbf{Generic task} & \textbf{Addressed in} \\
\midrule
\endfirsthead
\multicolumn{4}{@{}l}{\itshape Table~\ref{tab:crispdm-traceability} (continued)} \\
\toprule
\textbf{Phase} & \textbf{ID} & \textbf{Generic task} & \textbf{Addressed in} \\
\midrule
\endhead
\bottomrule
\endlastfoot
Business understanding & BU.1 & Determine business objectives & \crisploc{BU.1} \\
                       & BU.2 & Assess situation & \crisploc{BU.2} \\
                       & BU.3 & Determine data mining goals & \crisploc{BU.3} \\
                       & BU.4 & Produce project plan & \crisploc{BU.4} \\
\midrule
Data understanding     & DU.1 & Collect initial data & \crisploc{DU.1} \\
                       & DU.2 & Describe data & \crisploc{DU.2} \\
                       & DU.3 & Explore data & \crisploc{DU.3} \\
                       & DU.4 & Verify data quality & \crisploc{DU.4} \\
\midrule
Data preparation       & DP.1 & Select data & \crisploc{DP.1} \\
                       & DP.2 & Clean data & \crisploc{DP.2} \\
                       & DP.3 & Construct data & \crisploc{DP.3} \\
                       & DP.4 & Integrate data & \crisploc{DP.4} \\
                       & DP.5 & Format data & \crisploc{DP.5} \\
\midrule
Modeling               & MO.1 & Select modeling technique & \crisploc{MO.1} \\
                       & MO.2 & Generate test design & \crisploc{MO.2} \\
                       & MO.3 & Build model & \crisploc{MO.3} \\
                       & MO.4 & Assess model & \crisploc{MO.4} \\
\midrule
Evaluation             & EV.1 & Evaluate results & \crisploc{EV.1} \\
                       & EV.2 & Review process & \crisploc{EV.2} \\
                       & EV.3 & Determine next steps & \crisploc{EV.3} \\
\midrule
Deployment             & DE.1 & Plan deployment & \crisploc{DE.1} \\
                       & DE.2 & Plan monitoring and maintenance & \crisploc{DE.2} \\
                       & DE.3 & Produce final report &
  \hyperref[crisp:DE.3]{General Introduction}, p.~\pageref{crisp:DE.3} \\
                       & DE.4 & Review project &
  \hyperref[crisp:DE.4]{General Conclusion}, p.~\pageref{crisp:DE.4} \\
\end{longtable}
\endgroup


\section{System Architecture and Workflow}
% A high-level box-and-arrow diagram of the real pipeline (data sources ->
% preparation -> model -> outputs/integration) BEFORE the details below.

\section{Data Understanding}
% CRISP-DM phase 2. Numbers must be computed from the actual data.

\subsection{Initial Data Collection}
\crispdm{DU.1}
% DU.1 Sources, acquisition method and date, licence or permission, provenance,
% problems met during collection.

\subsection{Data Description}
\crispdm{DU.2}
% DU.2 Format, volume, schema/fields with types and units, classes and counts,
% time span. A table is usually clearest.

\subsection{Data Exploration}
\crispdm{DU.3}
% DU.3 Distributions, class balance, correlations, representative (anonymized)
% samples, and the findings that shaped later decisions.

\subsection{Data Quality Verification}
\crispdm{DU.4}
% DU.4 Missing values, label noise, duplicates, outliers, inconsistencies and
% bias, quantified. Every issue listed here is handled in Data Cleaning.

\section{Data Preparation}
% CRISP-DM phase 3.

\subsection{Data Selection}
\crispdm{DP.1}
% DP.1 What was kept or excluded (records, attributes, time ranges) and why.

\subsection{Data Cleaning}
\crispdm{DP.2}
% DP.2 For each quality issue from the previous phase: action taken and its
% effect (counts before/after), or why it was left alone.

\subsection{Data Construction and Feature Engineering}
\crispdm{DP.3}
% DP.3 Derived attributes (with formulas), augmentation or generated records
% (technique, parameters), and the reason for each.

\subsection{Data Integration}
\crispdm{DP.4}
% DP.4 How sources were merged (keys, conflicts). With a single source, change
% the marker to \crispdm[na]{DP.4} and explain in one or two sentences.

\subsection{Data Formatting}
\crispdm{DP.5}
% DP.5 Format conversions required by the modeling tools (label format,
% resizing, normalization, encoding, file layout) and a short description of
% the final dataset handed to modeling.

\section{Modeling}
% CRISP-DM phase 4. The test design (MO.2) and the model assessment (MO.4) are
% reported in Chapter 4; point to them here.

\subsection{Selection of Modeling Techniques}
\crispdm{MO.1}
% MO.1 Candidate techniques, the choice and its reasons (link to
% Chapter~\ref{chap:sota}), and the assumptions the technique makes.

\subsection{Model Building}
\crispdm{MO.3}
% MO.3 Architecture, hyperparameter table, training procedure (epochs,
% optimizer, schedule, seeds), hardware, training time, and training curves
% drawn from the real logs.

\section{Deployment and Integration}
% CRISP-DM phase 6 (planning). Keep "implemented", "prototype" and "planned"
% distinct. Never claim production use that was not demonstrated.

\subsection{Deployment Plan}
\crispdm{DE.1}
% DE.1 How the model reaches its users (API, application, edge device, batch
% job): what is implemented versus proposed; an integration diagram if useful.

\subsection{Monitoring and Maintenance Plan}
\crispdm{DE.2}
% DE.2 What would be monitored (input drift, prediction quality, latency),
% retraining triggers and ownership. Usually a proposal: label it as such.

\section{Conclusion}
% Summarize the methodology; transition to the results.
